% RESEARCH_DRAFT_SHA256: 79be705294255665a615bd90b95479a63ef62d5ea5337d75ffcbe302be9a42c2
% Layout and source-reference wrappers expanded for the web renderer.
\documentclass[11pt]{article}
\usepackage[margin=25mm]{geometry}
\usepackage{fontspec,unicode-math}
\setmainfont{Latin Modern Roman}
\setmathfont{Latin Modern Math}
\usepackage{amsmath,amssymb,mathtools,xurl,hyperref}
\hypersetup{colorlinks=true,urlcolor=blue,linkcolor=blue}
\setcounter{secnumdepth}{0}
\setlength{\parindent}{0pt}
\setlength{\parskip}{5pt}
\setlength{\emergencystretch}{3em}
\newcommand{\sref}[2]{\hyperlink{source:#1:#2}{[S#2]}}
\newcommand{\cataloguescope}[1]{\paragraph{Source record.}#1}
\newcommand{\floor}[1]{\left\lfloor#1\right\rfloor}
\begin{document}
\section*{Research-attempt provenance}
This standalone research attempt retains the complete original definition body below,
including its historical source and preparation notes. Those notes describe the
earlier catalogue preparation; the new attempt and its verification are identified
separately after that body. No original ID, catalogue status, or source statement is changed.

\section{Version-pinned workspace semantics for activities}\label{statement}

\subsection{Definitions and mathematical statement}

\paragraph{Shared notation.}Standard number systems are $\mathbb N,\mathbb Z,\mathbb Q,\mathbb R,\mathbb C$. Specialized notation and index ranges follow the cited source and the definitions below.


In an XML-based document model, an element represents a structural node and an attribute attaches a named value to that node. Workspace denotes blank area reserved for writing in an exercise or activity. Whether an attribute is allowed is a schema and implementation question.
\paragraph{Record scope and qualifications.}
The catalogue title is retained as metadata. Named certificates, contracts, bounds or frameworks in that title are not asserted as results of the original source extract.
The registry marks reconstructed wording. The statement's currently is historical and does not establish current schema support; the literal attribute spelling is preserved.

Catalogue screening limitation: PreTeXt software design, documentation, or publication feature request; no mathematical problem is specified by the supplied passage.

The source snapshot labels this record partially solved; this is catalogue metadata rather than independent verification of a mathematical result or the current completion of a project.
The source text below is rendered literally where necessary; damaged formulas and composite extracts are not silently repaired into a new mathematical claim.
The supplied catalogue marks this formulation as an unverified reconstruction; missing assumptions are not certified by the source link.
\paragraph{Statement recorded in the supplied source.}
Currently the activities element does not allow `@workspace`, maybe it should? \href{https://github.com/PreTeXtBook/community/wiki/Moderated-Problem-Session}{[S2]}
\paragraph{Source status.}
The supplied catalogue labels this source record partially solved; that metadata was checked against the supplied snapshot on 2026-10-11. The label is preserved for provenance and is not an independent verification of present open status or mathematical correctness.
\subsection{Short English statement}
Consider allowing a workspace attribute on an activity element.
\subsection{Sources}
\begin{itemize}
\item[S1] Ulam AI and the UnsolvedMath Contributors, UnsolvedMath: A Curated Collection of Open Mathematics Problems, catalogue entry. Catalogue curation and metadata: CC BY 4.0. The source extract is retained; definitions, English summaries and qualifications are editorial additions.
\url{https://huggingface.co/datasets/ulamai/UnsolvedMath}.
\item[S2] AIM Workshop Problem Lists, source linked by the catalogue.
\url{https://github.com/PreTeXtBook/community/wiki/Moderated-Problem-Session}.
\item[S3] Official PreTeXt Guide, Worksheets (context for worksheet and workspace terminology).
\url{https://pretextbook.org/doc/guide/html/overview-worksheet.html}.
\end{itemize}
\paragraph{Preparation and review.}
The original source and UnsolvedMath contributors retain attribution for the source extract. Codex (GPT-6; exact serving revision not exposed) assisted with definitions, English summaries, scope notes and formatting. No new human mathematical authorship or independently verified solution is asserted. Catalogue identity checked against the complete supplied registry and current repository. Distinct upstream records, including repeated questions and broad or damaged extracts, are retained for complete coverage of the source CS category. A repeated record is not a new independent conjecture. Complete mathematical-equivalence and current-literature review has not been performed. Preparation and final review follow contribution prompts C and D, with source gaps disclosed. Structural and rendering checks do not establish a new theorem or open status.
\label{end:statement}

\section{Literature and context check}
The community-session source was consulted directly. Its entry is a feature or design request, without a complete formal specification. The catalogue title is editorial and is not evidence of an existing theorem. The following finite model is an auxiliary investigation; no PreTeXt implementation or project completion is asserted. The consulted schema's displayed ProjectLike definition permits an optional workspace attribute and is used by the singular activity element. This is a reading of the displayed master snapshot, not a pinned release or a proof about a distinct activities element. The official worksheet guide describes print workspace and its suppression in HTML.
\section{Research attempt: Aggregate workspace versus inherited workspace}
To investigate the design ambiguity, let an activity have $n\ge1$ child tasks and a nonnegative workspace value $w$. Two plausible semantics are: each child inherits $w$, giving total blank height $nw$; or the activity has a total budget $w$, allocated as heights $h_i\ge0$ with $\sum_i h_i=w$. These disagree already for $n=2,w=20$: total height is $40$ under inheritance and $20$ under equal allocation. Merely accepting the same attribute in a grammar cannot determine which layout is intended.

Suppose each task requires a minimum blank height $a_i\ge0$. An allocation satisfying both the minima and the total budget exists exactly when $w\ge\sum_i a_i$. Necessity follows by summing $h_i\ge a_i$. For sufficiency, put $h_i=a_i$ for $i>1$ and $h_1=a_1+w-\sum_i a_i$. All heights are nonnegative, meet the minima, and sum to $w$. Equal allocation $h_i=w/n$ is not sufficient whenever a single $a_i>w/n$, even if a feasible allocation exists. For example $a=(18,2)$ and $w=20$ have feasible allocation $(18,2)$, while $(10,10)$ fails.

There is also a nesting obstruction. If an inherited value is expanded into physical space independently at each ancestor and at each leaf, a two-level activity can count the same intended budget more than once. An exact-budget interpretation instead associates each physical gap with one owner and sums it once. The finite sum argument then conserves the declared budget; it does not choose appropriate pedagogical heights or account for page breaks. These examples show why schema admission, inheritance, and rendering must be specified separately.
\section{Verification and adversarial audit}
The argument is an exact analysis of the stated finite model. Its counterexample was checked by direct substitution; no PreTeXt build, browser test, or user study was run. The model's assumptions are explicit and are not asserted to describe every real document.
\paragraph{Audit corrections.} The historical absence claim is not repeated as a current fact. Singular activity evidence is not generalized to an unverified plural element, and workspace semantics are selected explicitly.
\paragraph{External computation.} None; symbolic arguments and boundary checks only.
\section{Scope and remaining gaps}
The calculation identifies a concrete failed route or a restricted correctness property. Actual schema integration, transformations, usability, and compatibility with the full original feature request remain unimplemented and unverified. Thus the original agenda remains unresolved, with no claimed mathematical novelty.
\paragraph{Integration review.} Codex (GPT-6; exact serving revision not exposed) reviewed the written arguments and expanded layout and source-reference wrappers. Repository metadata and the production web renderer were checked using Python 3.13 and JavaScript with MathJax 3.2.2. This model-assisted review is not independent human proof verification; no PDF compilation or formal proof checking is asserted.

\section{Completion Estimate}
\noindent\textbf{COMPLETION: 0\%}
The calculations below the preserved source statement investigate a concrete route,
but establish no substantiated new progress toward the entire original target.
Known special cases and failed sufficient conditions are not credited as a solution.
\section{References}
The original source references and catalogue attribution are preserved in the
Sources section above. The following URLs were actually consulted for this attempt;
retention of an additional original reference does not claim fresh consultation of it.
\begin{itemize}
\item \url{https://github.com/PreTeXtBook/community/wiki/Moderated-Problem-Session} (consulted 2026-10-11). 
\item \url{https://pretextbook.org/doc/guide/html/overview-worksheet.html} (consulted 2026-10-11). 
\item \url{https://github.com/PreTeXtBook/pretext/blob/master/schema/pretext.rnc} (consulted 2026-10-11). 
\end{itemize}
\section{Model and contributor disclosure}
OpenAI gpt-6-astra, requested reasoning configuration ultra, generated this research
attempt on 2026-10-11. The exact serving revision is not exposed. Codex GPT-6
performed orchestration, repository-format review, and a model-assisted mathematical
audit of the written arguments; its exact serving revision
is also not exposed. The mathematical route, boundary checks and final adversarial
self-audit followed contribution prompts A then B in
\texttt{docs/contribution\_prompts.md}, after reading \texttt{CONTRIBUTING.md}.
Python 3.14.2 assembled the standalone documents; any mathematical computation is
specified in Verification. No theorem prover was used. No human mathematical
authorship, substantial human proof correction, or independent human proof checking
is asserted. Model-assisted review, source retrieval, and structural checks do not
establish correctness or novelty. Full private reasoning transcripts are unavailable
as reproducibility artifacts; the complete arguments, inputs and outcomes actually
used are stated here or in the companion batch audit.
\section{Final status}
Unresolved as a new research contribution. This is a concrete route investigation,
not a claim of new complete resolution or a current-status certification. A known
answer reproduced for context does not become a new contribution. The catalogue's
source-status label is preserved; this attempt status does not assert that a
source-labelled solved question is open.
\end{document}
